\chapter{Équations algébriques non linéaires}

\section{Introduction}

Les équations algébriques non linéaires se présentent sous la forme
$f(x) = 0$ ou $f_i(x) = 0$ avec $i$ allant de $1$ à $n$. La variable $x$ peut
être réelle ou complexe. La résolution numérique de ces équations fait appel à
plusieurs méthodes qui seront expliquées dans ce chapitre. Il s'agit notamment
de la méthode de Newton-Raphson, la méthode de dichotomie et de la méthode de
Bairstow. Ce chapitre expliquera aussi comment effectuer les opérations sur
les polynômes ainsi que la présentation des méthodes permettant de trouver
plusieurs racines simultanément. Le cas des systèmes d'équations non linéaires
sera aussi présenté.

\section{Méthodes de Newton-Raphson et dichotomie}

La formule de Newton-Raphson a été donnée par Raphson en 1698. Mais Newton
avait suggéré une méthode plus voisine quelques mois auparavant. Il se
pourrait aussi que cette formule ait été utilisée par Héron en l'an 100 avant
Jésus-Christ pour trouver la racine carrée d'un nombre. Cette méthode est de
la classe des méthodes itératives par approximations successives qui se
traduisent par répétition systématique d'une opération déterminée, ce qui
finit par donner de proche en proche la solution exacte.

\subsection{Formule de la méthode de Newton-Raphson}

On commence généralement par trouver une solution grossière ou approchée
$x_0$, dite approximation d'ordre zéro. Dans le cas d'un problème de physique,
une telle solution $x_0$ peut être suggérée par la signification du problème
ou par une courbe. Généralement, on localise un intervalle $[a,b]$ tel que
$x_0$ appartienne à cet intervalle. Si nous désignons la solution exacte par
$\bar{x} = x_0 + h$, alors, la formule de Taylor donne
\[
  f(x_0 + h) \approx f(x_0) + h f'(x_0) + O(h^2)
\]
or
\[
  f(x_0 + h) = 0 \;\longrightarrow\; f(x_0) + h f'(x_0) \approx 0
  \;\longrightarrow\; h \approx -\frac{f(x_0)}{f'(x_0)}
\]
d'où
\[
  \bar{x} = x_0 - \frac{f(x_0)}{f'(x_0)} .
\]
En posant $x_1 = x_0 - \dfrac{f(x_0)}{f'(x_0)}$, on obtient une première
approximation de la solution (l'indice désignant l'ordre de l'approximation).
Procédons de même avec $x_1$, on obtient la seconde approximation,
\[
  x_2 = x_1 - \frac{f(x_1)}{f'(x_1)} .
\]
À l'ordre $n+1$, on obtient
\[
  x_{n+1} = x_n - \frac{f(x_n)}{f'(x_n)} .
\]
C'est une formule résultant d'une erreur d'ordre $O(h^2)$. En négligeant les
termes d'ordre supérieur, on a remplacé la courbe représentative de $f(x)$ par
une tangente en $x = x_0$ ; la méthode consiste géométriquement à tracer de
proche en proche les tangentes à la courbe et à rechercher leur point
d'intersection avec l'axe des abscisses. En effet, on a l'équation
\[
  f(\bar{x}) = y = f(x_0) + h f'(x_0) = f(x_0) + (\bar{x} - x_0) f'(x_0)
\]
qui est l'équation de la tangente au point $\bigl(x_0, f(x_0)\bigr)$.

Cette interprétation graphique fait que la méthode de Newton est souvent
appelée \emph{méthode de la tangente}.

\paragraph{Remarques} Lorsque $f(x_0)f''(x_0) > 0$, la suite
$x_{n+1} = x_n - \dfrac{f(x_n)}{f'(x_n)}$ est monotone, bornée et converge vers
$\bar{x}$. Ce résultat donne une condition suffisante pour que la méthode de
Newton converge vers $\bar{x}$. Mais cette condition n'est pas nécessaire.
Dans le cas où il n'y a pas convergence, on change la valeur de $x_0$.

\subsection{Condition d'arrêt}

La méthode de Newton converge vers la solution exacte $\bar{x}$, mais on ne
connaît pas explicitement à quel moment une certaine précision est atteinte.
Pour cela, on utilise généralement le résultat suivant :

Soit $\bar{x}$ une racine de l'équation $f(x) = 0$ et soit $x_n$ une valeur
approchée de $\bar{x}$. Si sur $[a,b]$ contenant $\bar{x}$ et $x_n$, on a
\[
  \abs{f'(x)} \ge m > 0, \quad \text{alors} \quad
  \abs{x_n - \bar{x}} \le \frac{\abs{f(x_n)}}{m} .
\]
Dans la pratique, on pourra prendre $m$ comme
$m = \min\bigl(\abs{f'(a)}, \abs{f'(b)}\bigr)$.

Si $E$ est la précision exigée, on s'arrêtera lorsque
$\dfrac{\abs{f(x_n)}}{m} < E$.

\begin{exercice}
Écrire l'algorithme de la méthode de Newton-Raphson en y intégrant la
condition d'arrêt.
\end{exercice}

\subsection{Autres formules : formule du troisième et du premier ordre}

La formule de Newton est l'une des méthodes les plus utilisées parce qu'elle
est relativement simple et d'une convergence très suffisante dans la plupart
des cas. Il faut toutefois que $x_0$ soit assez près de la racine pour que la
dérivée $f'$ ne change pas de signe entre $x_0$ et $x$. Si en effet $f'$
s'annule sans que $f$ s'annule aussi, la tangente à la courbe va couper l'axe
des $x$ à l'infini. On risque alors d'osciller indéfiniment autour de la
solution sans jamais l'atteindre.

Il est possible de modifier la formule de Newton et d'écrire
\[
  x_{n+1} = x_n - p\,\frac{f(x_n)}{f'(x_n)}
\]
où $p$ est un nombre ou une fonction. On peut aussi utiliser la formule
\[
  x_{n+1} = x_n - \frac{g\,f(x_n)}{g\,f'(x_n) + h\,f(x_n)} .
\]
En développant la série de Taylor jusqu'à l'ordre 2, on obtient
\[
  f(x_0 + h) = f(x_0) + h f'(x_0) + \frac{h^2}{2} f''(x_0) + O(h^3) .
\]
En remplaçant $h$ par $h = \bar{x} - x_0 \approx -\dfrac{f}{f'}$ (d'après la
formule de Newton), on obtient
\[
  x_{n+1} = x_n - \frac{f(x_n)}{f'(x_n)}
          - \frac{1}{2}\,\frac{f''(x_n)}{f'(x_n)}
            \left(\frac{f(x_n)}{f'(x_n)}\right)^{2} .
\]
C'est la formule du troisième ordre.

Soit l'équation $f(x) = 0$. De cette équation, on peut tirer $x$ en fonction
de $x$ en écrivant $x = g(x)$, où $g$ est défini à partir de $f(x)$. Dans ce
cas on peut écrire $x_{n+1} = g(x_n)$. C'est la formule d'itération du premier
ordre.

\subsection{Formule de Newton pour les racines complexes}

Soit l'équation
\[
  f(z) = 0
\]
avec $z = x + iy$. On a
\[
  f(z) \approx f(z_0) + (z - z_0) f'(z_0)
  \;\Longrightarrow\; z - z_0 = -\frac{f(z_0)}{f'(z_0)} .
\]
Pour le calcul numérique, il faut séparer les parties réelles et imaginaires.
On aura donc
\[
  f(z_0) = G(x_0, y_0) + iH(x_0, y_0), \qquad
  \frac{\dd f}{\dd z_0} = f'(z_0) = G'(x_0, y_0) + iH'(x_0, y_0)
\]
et
\[
  \left\{
  \begin{aligned}
    x_{n+1} &= x_n - \left.\frac{GG' + HH'}{G'^2 + H'^2}\right|_{z_n}\\[0.4em]
    y_{n+1} &= y_n - \left.\frac{HG' - GH'}{G'^2 + H'^2}\right|_{z_n}
  \end{aligned}
  \right.
\]

\subsection{Méthode de dichotomie ou de bissection}

Supposons que $f(a) < 0 < f(b)$ et $f$ est strictement croissante sur $[a,b]$.
La solution $x_0$ de l'équation $f(x) = 0$ appartient à l'intervalle $[a,b]$.
Le milieu de l'intervalle est $c = (a+b)/2$. Si $f(a)$ et $f(c)$ ont des signes
opposés, alors $x_0$ appartient à $[a,c]$. Si par contre $f(b)$ et $f(c)$ sont
de signes opposés, alors $x_0$ appartient à $[c,b]$. Dans l'un ou l'autre
cas, la taille de l'intervalle est divisée par 2. Une nouvelle bissection peut
être faite sur le nouvel intervalle contenant la solution et ainsi de suite.
Ainsi, après $n$ étapes, la taille de l'intervalle devient $(b-a)/2^n$.

\begin{exercice}
Écrire l'algorithme de la méthode de bissection.
\end{exercice}

\section{Calcul de toutes les racines d'une équation polynomiale}

\subsection{Opérations sur les polynômes}

\subsubsection{Schéma de calcul numérique d'un polynôme}

Pour déterminer les valeurs numériques de
\[
  P_n(x) = a_0x^n + a_1x^{n-1} + \dots + a_{n-1}x + a_n ,
\]
on peut calculer chaque terme et faire la somme. Mais, ce procédé n'est pas en
général court. Il exige $2n-1$ multiplications. On peut utiliser le schéma dit
de \textbf{Horner}. On calcule les nombres $b_k$ définis par :
\[
  \begin{aligned}
    b_0 &= a_0,\\
    b_k &= a_k + b_{k-1}x \qquad (k \text{ allant de } 1 \text{ à } n).
  \end{aligned}
\]
$b_n$ est alors la valeur numérique cherchée.

L'algorithme utilisé se présente ainsi :

\begin{algo}
  \State \textbf{Début}
  \State \quad $x \leftarrow$ \textit{valeur de} $x$
  \State \quad $b_0 \leftarrow a_0$
  \State \quad \textbf{Pour} $k$ allant de $1$ à $n$ \textbf{faire}
  \State \quad\quad $b_k \leftarrow a_k + b_{k-1}x$
  \State \quad \textbf{Fin pour}
  \State \quad \textbf{Écrire} $b_n = P_n(x)$
  \State \textbf{Fin}
\end{algo}

La première méthode de calcul exige $(2n-1)$ multiplications alors que le
schéma de Horner n'en demande que $n$ multiplications et le nombre d'additions
est le même dans les deux cas. De plus le schéma de Horner a l'avantage de
n'introduire qu'un type d'opérations :
\[
  b_k = a_k + b_{k-1}x .
\]
La programmation est alors plus aisée.

Pour un polynôme à coefficients complexes,
\[
  \begin{aligned}
    &P_n(z) = c_0z^n + c_1z^{n-1} + \dots + c_{n-1}z + c_n,\\
    &z = x + iy \quad\text{et}\quad c_p = a_p + ib_p,\\
    &P_n(z) = G_n(x,y) + iH_n(x,y).
  \end{aligned}
\]
Le schéma de Horner peut également être utilisé. On remplace tout simplement
les quantités $b_k$ par les quantités
\[
  P_k = G_k + iH_k
\]
avec
\[
  \begin{aligned}
    G_0 &= a_0, & H_0 &= b_0,\\
    G_k &= a_k + xG_{k-1} - yH_{k-1}, &
    H_k &= b_k + yG_{k-1} + xH_{k-1}.
  \end{aligned}
\]

\subsubsection{Division d'un polynôme par \texorpdfstring{$x-r$}{x-r}}

En divisant le polynôme
\[
  P_n(x) = a_0x^n + a_1x^{n-1} + \dots + a_{n-1}x + a_n
\]
par
\[
  C(x) = x - r ,
\]
on obtient un quotient $Q(x)$ et un reste $R$ liés entre eux par la relation
\[
  P(x) = (x - r)\,Q(x) + R .
\]
$Q(x)$ est un polynôme de degré $n-1$ et $R$ est une constante. Ils sont de la
forme
\[
  \begin{aligned}
    Q(x) &= b_0x^{n-1} + b_1x^{n-2} + \dots + b_{n-2}x + b_{n-1}\\
    R &= b_n
  \end{aligned}
\]
où $b_k = a_k + r\,b_{k-1}$ $(k = 1 \text{ à } n)$, avec $b_0 = a_0$.

Si $b_n = 0$, alors le polynôme est divisible par $x - r$ ; c'est-à-dire que
$r$ est une racine de $P(x)$.

\paragraph{Remarques}
\begin{itemize}
  \item[$*$] Ces formules sont semblables pour des grandeurs complexes. La
        seule différence apparaît dans le calcul effectif qui exige un
        dédoublement de colonne comme ci-dessus.
  \item[$**$] Quand les coefficients $a_k$ sont réels et que
        $r = \alpha + i\beta$, on peut diviser $P(x)$ par le facteur
        quadratique $(x - r)(x - r^*)$ où $r^*$ est le complexe conjugué.
\end{itemize}

\subsubsection{Division d'un polynôme par un facteur quadratique}

Soit le polynôme $P_n(x)$ que l'on désire diviser par le facteur quadratique
\[
  Q_2(x) = x^2 - Sx + P .
\]
On a :
\[
  P_n(x) = Q_2(x)\,Q_{n-2}(x) + R_1(x) .
\]
$Q_{n-2}(x)$ et $R_1(x)$ sont respectivement les polynômes de degré $n-2$ et
$1$ que l'on peut écrire sous la forme
\[
  \begin{aligned}
    Q_{n-2}(x) &= b_0x^{n-2} + b_1x^{n-3} + \dots + b_{n-3}x + b_{n-2}\\
    R_1(x) &= b_{n-1}(x - S) + b_n
  \end{aligned}
\]
avec
\[
  \begin{aligned}
    b_0 &= a_0, \qquad b_1 = a_1 + Sb_0,\\
    b_2 &= a_2 + Sb_1 - Pb_0,\\
    &\;\;\vdots\\
    b_k &= a_k + Sb_{k-1} - Pb_{k-2} \qquad (k = 2 \text{ à } n).
  \end{aligned}
\]
Si $S$ et $P$ sont respectivement la somme et le produit de deux racines du
polynôme $P_n(x)$, alors le reste $R_1$ est nul ; c'est-à-dire $b_{n-1} = 0$
et $b_n = 0$.

\subsection{Calcul d'une racine réelle et passage à l'équation de degré
\texorpdfstring{$n-1$}{n-1}}

L'idée la plus simple est de calculer la suite des valeurs
$b_k = a_k + b_{k-1}x_0$ qui conduisent à la valeur du polynôme et comme on
devrait avoir $b_n = 0$, on obtiendra
\[
  x = -\frac{a_n}{b_{n-1}} .
\]
On recommence le calcul des quantités $b_k$ avec cette nouvelle valeur de $x$
jusqu'à obtenir la précision choisie. À ce moment les coefficients $b_0$,
$b_1$, \dots, $b_{n-1}$ donnent le polynôme restant de degré $n-1$. On peut
alors continuer le processus de recherche des racines sur le polynôme de
degré $n-1$.
\[
  P_n(x) = a_0x^n + a_1x^{n-1} + \dots + a_{n-1}x + a_n
\]
L'inconvénient de ce processus est qu'il n'est pas toujours convergent. Cette
méthode appelée méthode de \textbf{Lin} fut développée en 1941.

Elle peut également s'appliquer sur le plan complexe. Pour cela, on part d'une
valeur arbitraire $z_0 = x_0 + iy_0$ et on calcule les coefficients de Horner
complexes jusqu'à $G_{n-1}$ et $H_{n-1}$. On obtient alors des nouvelles
valeurs $z_1 = x_1 + iy_1$ en écrivant que $G_n = 0$ et $H_n = 0$. On a en
effet
\[
  \left\{
  \begin{aligned}
    a_n + x_1G_{n-1} - y_1H_{n-1} &= 0\\
    b_n + x_1H_{n-1} + y_1G_{n-1} &= 0
  \end{aligned}
  \right.
  \quad\Longrightarrow\quad
  \left\{
  \begin{aligned}
    x_1 &= -\frac{a_nG_{n-1} + b_nH_{n-1}}{G_{n-1}^2 + H_{n-1}^2}\\[0.4em]
    y_1 &= \frac{a_nH_{n-1} - b_nG_{n-1}}{G_{n-1}^2 + H_{n-1}^2}
  \end{aligned}
  \right.
\]
On recommence le processus jusqu'à ce que
\[
  \abs{z^{(k+1)} - z^{(k)}} < \varepsilon .
\]

\begin{exercice}
Établir l'algorithme de ce schéma de Lin pour les solutions réelles.
\end{exercice}

\subsection{Formule d'itération de Newton sur l'axe réel}

Soit l'équation $P_n(x) = 0$, la dérivée de ce polynôme en $x$ est $P'_n(x)$.
Les valeurs numériques de $P_n(x)$ et de $P'_n(x)$ se calculent par le schéma
de Horner suivant :
\[
  \begin{aligned}
    b_0 &= a_0\\
    b_k &= a_k + xb_{k-1} \qquad (k = 1, \dots, n)
  \end{aligned}
\]
et
\[
  \begin{aligned}
    c_0 &= b_0\\
    c_k &= b_k + xc_{k-1} \qquad (k = 1, \dots, n-1)
  \end{aligned}
\]
On a
\[
  \left\{
  \begin{aligned}
    P_n(x) &= b_n\\
    P'_n(x) &= c_{n-1}
  \end{aligned}
  \right.
\]
La formule d'itération de Newton devient
\[
  x_{k+1} = f(x_k) \qquad\text{où}\qquad f(x) = x - \frac{b_n}{c_{n-1}} .
\]
On peut alors poursuivre le processus jusqu'à obtenir
\[
  \abs{x_{k+1} - x_k} < \varepsilon .
\]
Après avoir trouvé une racine $x^*$ de l'équation, on divise le polynôme par
$x - x^*$ et on continue le processus jusqu'à l'obtention de toutes les
racines.

\subsection{Calcul de deux racines et passage à l'équation de degré
\texorpdfstring{$n-2$}{n-2}}

Dans le cas de la division d'un polynôme par le facteur quadratique
\[
  G_2(x) = x^2 - Sx + P ,
\]
si $S$ et $P$ sont la somme et le produit de deux racines du polynôme
$P_n(x)$, alors on aura
\[
  b_n = 0 \quad\text{et}\quad b_{n-1} = 0
  \quad\Longrightarrow\quad
  P = \frac{a_n}{b_{n-2}} \quad\text{et}\quad
  S = \frac{Pb_{n-3} - a_{n-1}}{b_{n-2}} .
\]
Le processus peut alors continuer jusqu'à $b_n = b_{n-1} = 0$ à la précision
choisie. On obtient finalement un nouveau polynôme de degré $n-2$ pour lequel
on peut reprendre le calcul pour retrouver les deux racines suivantes et ainsi
de suite, on parvient à résoudre l'équation.

\subsection{Méthode de Bairstow}

Soit le polynôme :
\[
  P_n(x) = a_0x^n + a_1x^{n-1} + \dots + a_{n-1}x + a_n
\]
à diviser par $G_2(x) = x^2 - Sx + P$. On obtient
\[
  P_n(x) = G_2(x)\,Q_{n-2}(x) + R_1(x)
\]
où
\[
  \begin{aligned}
    Q_{n-2}(x) &= b_0x^{n-2} + b_1x^{n-3} + \dots + b_{n-3}x + b_{n-2}\\
    \text{et}\quad R_1(x) &= b_{n-1}(x - S) + b_n
  \end{aligned}
\]
avec
\[
  b_0 = a_0, \qquad b_1 = a_1 + Sb_0, \qquad
  b_k = a_k + Sb_{k-1} - Pb_{k-2} \quad (k = 2, \dots, n).
\]
Si l'on a $b_n = 0$ et $b_{n-1} = 0$, alors $S$ et $P$ sont la somme et le
produit de deux racines du polynôme et les grandeurs $b_k$ sont les
coefficients du polynôme de degré $n-2$. La méthode de Bairstow consiste à
partir des valeurs arbitraires de $S$ et $P$ et de calculer la suite des
coefficients $b_k$.

Soit
\[
  \begin{aligned}
    F(S,P) &= b_{n-1},\\
    G(S,P) &= b_n .
  \end{aligned}
\]
Si $F$ et $G$ ne sont pas nuls à la précision choisie, alors on améliore $S_0$
et $P_0$ ($S_0$ et $P_0$ étant les valeurs initiales de $S$ et de $P$) par la
formule d'itération de Newton à deux variables. Les deux fonctions de $S$ et
de $P$ sont $F$ et $G$ et les formules d'itération se mettent sous la forme
\[
  \begin{aligned}
    S_{k+1} &= S_k + \frac{\alpha}{\Delta}\\
    P_{k+1} &= P_k + \frac{\beta}{\Delta}
  \end{aligned}
\]
avec
\[
  \begin{aligned}
    \alpha &= F\frac{\partial G}{\partial P} - G\frac{\partial F}{\partial P}\\
    \beta &= G\frac{\partial F}{\partial S} - F\frac{\partial G}{\partial S}\\
    \Delta &= \frac{\partial G}{\partial S}\frac{\partial F}{\partial P}
            - \frac{\partial F}{\partial S}\frac{\partial G}{\partial P}
  \end{aligned}
\]
Pour obtenir les dérivées partielles par rapport à $S$, il suffit de dériver
les relations $b_k$ par rapport à $S$. Pour cela nous posons
\[
  \frac{\partial b_k}{\partial S} = c_{k-1} .
\]
En dérivant $b_k = a_k + Sb_{k-1} - Pb_{k-2}$ par rapport à $S$, on a
\[
  \frac{\partial b_k}{\partial S}
  = b_{k-1} + S\frac{\partial b_{k-1}}{\partial S} - P\frac{\partial b_{k-2}}{\partial S},
  \qquad\text{soit}\qquad
  c_{k-1} = b_{k-1} + Sc_{k-2} - Pc_{k-3} .
\]
On obtient alors la suite suivante
\[
  \begin{aligned}
    c_0 &= b_0\\
    c_1 &= b_1 + Sc_0\\
    c_k &= b_k + Sc_{k-1} - Pc_{k-2} \qquad (k = 2, \dots, n)
  \end{aligned}
\]
Posons
\[
  \frac{\partial b_k}{\partial P} = -c_{k-2} .
\]
On constate que l'on obtient une suite identique à celle définie ci-dessus.

Il vient finalement :
\[
  \begin{aligned}
    \alpha &= b_nc_{n-3} - b_{n-1}c_{n-2}\\
    \beta &= b_nc_{n-2} - b_{n-1}c_{n-1}\\
    \Delta &= c_{n-2}^2 - c_{n-1}c_{n-3}
  \end{aligned}
\]
Ces formules ont été établies par Bairstow en 1914 pour des besoins de
l'industrie d'aéronautique. Avec ces nouvelles valeurs de $S$ et de $P$, on
calcule à nouveau les coefficients $b_q$ et $c_q$ ainsi de suite jusqu'à ce
que $b_{n-1}$ et $b_n$ soient nuls à la précision choisie. Une fois que l'on a
terminé le calcul de $S$ et de $P$, on détermine les deux racines
correspondantes par :
\[
  x = \frac{1}{2}\Bigl[S \pm \sqrt{S^2 - 4P}\Bigr]
  \qquad\text{ou}\qquad
  x + iy = \frac{1}{2}\Bigl[S \pm i\sqrt{-S^2 + 4P}\Bigr] .
\]
Les coefficients restants $b_k$ sont ceux de l'équation de degré $n-2$. On peut
alors recommencer le processus et calculer par paire toutes les racines.
Toutefois, si $n$ est impair, il restera à la fin une équation de la forme :
\[
  b_0 + b_1x = 0 \quad\Longrightarrow\quad x = -\frac{b_0}{b_1} .
\]

\subsection{Calcul simultané de toutes les racines}

Pour calculer simultanément toutes les racines $x_1, x_2, \dots, x_n$ d'un
polynôme de degré $n$, il faut disposer de $n$ équations entre ces racines.
Malheureusement, c'est souvent difficile d'obtenir ce système d'équations
entre les racines. Mais pour les équations de degré inférieur à 5, on peut
trouver des relations : relations de Viète.

\subsubsection{Cas où toutes les racines sont réelles}

Si nous considérons l'équation du 3\textsuperscript{e} degré
\[
  x^3 + a_1x^2 + a_2x + a_3 = 0 .
\]
On a les relations de Viète suivantes entre les racines :
\[
  \left\{
  \begin{aligned}
    x_1 + x_2 + x_3 &= -a_1\\
    x_1x_2 + x_1x_3 + x_2x_3 &= a_2\\
    x_1x_2x_3 &= -a_3
  \end{aligned}
  \right.
  \quad\Longleftrightarrow\quad
  \left\{
  \begin{aligned}
    x_1 &= -a_1 - x_2 - x_3\\
    x_2 &= \frac{1}{x_1}\bigl(a_2 - x_1x_3 - x_2x_3\bigr)\\
    x_3 &= -\frac{a_3}{x_1x_2}
  \end{aligned}
  \right.
\]
On obtient ainsi un schéma itératif du 1\textsuperscript{er} ordre qui permet
par approximations successives d'avoir les trois racines de l'équation.
\[
  \left\{
  \begin{aligned}
    x_1^{(k+1)} &= -a_1 - x_2^{(k)} - x_3^{(k)}\\
    x_2^{(k+1)} &= \frac{1}{x_1^{(k)}}\Bigl(a_2 - x_1^{(k)}x_3^{(k)} - x_2^{(k)}x_3^{(k)}\Bigr)\\
    x_3^{(k+1)} &= -\frac{a_3}{x_1^{(k)}x_2^{(k)}}
  \end{aligned}
  \right.
\]

\subsubsection{Cas où il y a des racines complexes}

Si les coefficients de l'équation sont réels, les racines complexes se
présentent par paires conjuguées l'une de l'autre. En associant les racines
par couple et en faisant apparaître les sommes et les produits, on opère
uniquement sur les quantités réelles.

Pour une équation du 3\textsuperscript{e} degré :
$x^3 + a_1x^2 + a_2x + a_3 = 0$, si on pose $x_1 + x_2 = S$, $x_1x_2 = P$,
$x_3 = r$, les relations de Viète s'écrivent :
\[
  \left\{
  \begin{aligned}
    S + r + a_1 &= 0\\
    P + rS - a_2 &= 0\\
    rP + a_3 &= 0
  \end{aligned}
  \right.
\]
De ces relations, on tire :
\[
  \left\{
  \begin{aligned}
    S &= -r - a_1\\
    P &= a_2 - rS\\
    r &= -\frac{a_3}{P}
  \end{aligned}
  \right.
\]
On peut alors partir des valeurs quelconques de $S$, $P$ et $r$ pour faire le
processus itératif et déterminer ces grandeurs à la précision choisie.

\subsubsection{Extension aux équations de degré supérieur}

Pour une équation du 4\textsuperscript{e} ordre :
$x^4 + a_1x^3 + a_2x^2 + a_3x + a_4 = 0$, on a les relations suivantes :
\[
  \left\{
  \begin{aligned}
    S_1 + S_2 + a_1 &= 0\\
    S_1S_2 + P_1 + P_2 - a_2 &= 0\\
    S_1P_2 + S_2P_1 + a_3 &= 0\\
    P_1P_2 - a_4 &= 0
  \end{aligned}
  \right.
\]
avec
\[
  S_1 = x_1 + x_2, \quad P_1 = x_1x_2, \quad S_2 = x_3 + x_4
  \quad\text{et}\quad P_2 = x_3x_4 .
\]
On peut alors tirer le schéma itératif du 1\textsuperscript{er} ordre
suivant :
\[
  \left\{
  \begin{aligned}
    S_1 &= -a_1 - S_2\\
    P_1 &= a_2 - P_2 - S_1S_2\\
    S_2 &= -\frac{1}{P_1}\bigl(a_3 + S_1P_2\bigr)\\
    P_2 &= \frac{a_4}{P_1}
  \end{aligned}
  \right.
\]
\[
  \begin{aligned}
    x_1 &= \frac{S_1 - \sqrt{S_1^2 - 4P_1}}{2}\,; &
    x_2 &= \frac{S_1 + \sqrt{S_1^2 - 4P_1}}{2}\\[0.4em]
    x_3 &= \frac{S_2 - \sqrt{S_2^2 - 4P_2}}{2}\,; &
    x_4 &= \frac{S_2 + \sqrt{S_2^2 - 4P_2}}{2}
  \end{aligned}
\]
Pour une équation du 5\textsuperscript{e} degré on a :
\[
  \left\{
  \begin{aligned}
    S_1 &= -\bigl(a_1 + S_2 + r\bigr)\\
    P_1 &= a_2 - P_2 - S_1S_2 - r\bigl(S_1 + S_2\bigr)\\
    S_2 &= -\frac{1}{P_1}\Bigl[a_3 + S_1P_2 + r\bigl(P_1 + P_2 + S_1S_2\bigr)\Bigr]\\
    P_2 &= \frac{1}{P_1}\Bigl[a_4 - r\bigl(S_1P_2 + S_2P_1\bigr)\Bigr]\\
    r &= -\frac{a_5}{P_1P_2}
  \end{aligned}
  \right.
  \qquad\text{où } r = x_5 .
\]

\subsubsection{Quelques relations et règles importantes}

\paragraph{Relations de Viète}
Soit le polynôme de degré $n$ en $z$ défini par :
\[
  P_n(z) = a_0z^n + a_1z^{n-1} + \dots + a_{n-1}z + a_n
\]
où les $a_i$ et $z$ sont des réels ou complexes.

Soit $z_j$ les racines de ce polynôme, on a alors
\[
  P_n(z) = a_0(z - z_1)(z - z_2)\cdots(z - z_n) .
\]
En développant cette expression et après identification des coefficients on
obtient les relations suivantes :
\[
  \begin{aligned}
    a_1 &= -a_0\sum_i z_i\\
    a_2 &= a_0\sum_{i<j} z_iz_j\\
    a_3 &= -a_0\sum_{i<j<k} z_iz_jz_k\\
    &\;\;\vdots\\
    a_n &= (-1)^n a_0\,z_1z_2\cdots z_n
  \end{aligned}
\]

\paragraph{Localisation des racines dans le plan}

\subparagraph{-- Règle des signes de Descartes.}
Le nombre de racines réelles positives ne peut pas dépasser le nombre de
changements de signe que l'on observe en parcourant la suite des coefficients
du polynôme. Il ne peut différer de ce nombre que par un nombre pair. On a une
règle analogue pour les racines négatives. Il suffit de changer $x$ en $-x$.

\textbf{Exemples :} L'équation $x^4 - 8x^3 + 24x^2 - 32x + 15 = 0$ n'a pas de
racines réelles négatives. De même, l'équation
$x^4 - 2x^3 - 3x^2 + 8x - 4 = 0$ n'a pas plus de trois racines réelles
positives.

\subparagraph{-- Règle de Gua.}
Si le carré d'un coefficient intermédiaire est inférieur ou égal au produit
des coefficients voisins, alors il y a des racines imaginaires.

\textbf{Exemple :} L'équation $x^3 + 7x^2 + x + 8 = 0$ possède une racine
imaginaire car on a
\[
  1 \times 1 < 7 \times 8 .
\]

\subparagraph{-- Règle des lacunes.}
S'il manque un terme entre deux termes de même signe ou plusieurs termes entre
deux termes de signe quelconque, alors il y a des racines imaginaires.

\textbf{Exemple :} L'équation $x^3 + 6x + 20 = 0$ a le terme en $x^2$ absent
et les coefficients des termes $x^3$ et $x$ sont positifs. On peut donc
affirmer qu'il y a des racines complexes.

\subparagraph{-- Règle de Sturm.}
Soit $a$ un nombre réel positif, si l'équation $(x - a)P_n(x) = 0$ présente
$(2k+1)$ variations de signe de plus que l'équation $P_n(x) = 0$ alors il y a
au moins $2k$ racines complexes.

\textbf{Exemple :} Soit l'équation $x^3 + 7x^2 + x + 7 = 0$ et l'équation
$(x - 1)P_n(x) = 0$ égale à
\[
  x^4 + 6x^3 - 6x^2 + 6x - 7 = 0
\]
qui présente trois changements de signe. Il y a donc 2 racines complexes pour
l'équation $P_n(x) = 0$.

\section{Système d'équations non linéaires}

Soit un système d'équations non linéaires
\[
  f_j(x_i) = 0 \qquad
  \left\{
  \begin{aligned}
    i &= 1, \dots, n\\
    j &= 1, \dots, n
  \end{aligned}
  \right.
\]
On peut aussi élaborer des processus itératifs du 1\textsuperscript{er},
2\textsuperscript{e} et 3\textsuperscript{e} ordre pour ce type de système
d'équations. En particulier, pour le système de deux équations non linéaires
\[
  \left\{
  \begin{aligned}
    f_1(x,y) &= 0\\
    f_2(x,y) &= 0
  \end{aligned}
  \right. ,
\]
on peut trouver $x$ à partir de $f_1(x,y)$ et $y$ à partir de $f_2(x,y)$ et
utiliser le processus itératif du 1\textsuperscript{er} ordre suivant :
\[
  \left\{
  \begin{aligned}
    x_{k+1} &= F_1(x_k, y_k)\\
    y_{k+1} &= F_2(x_k, y_k)
  \end{aligned}
  \right.
\]
On peut également généraliser les formules itératives de Jacobi et de
Gauss-Seidel. De même, à partir de la méthode de Newton, le système de deux
équations non linéaires à 2 inconnues donne les formules itératives
suivantes :
\[
  \begin{aligned}
    x_{k+1} &= x_k + \frac{1}{\Delta}\bigl(f_2f_{1,y} - f_1f_{2,y}\bigr)_{(x_k,y_k)}
      \qquad\text{(quantité évaluée au point } (x_k,y_k)\text{)}\\
    y_{k+1} &= y_k + \frac{1}{\Delta}\bigl(f_1f_{2,x} - f_2f_{1,x}\bigr)_{(x_k,y_k)}\\
    \Delta &= f_{1,x}f_{2,y} - f_{1,y}f_{2,x}, \qquad
    f_{1,x} = \frac{\partial f_1}{\partial x}, \;\dots
  \end{aligned}
\]
Dans le cas d'un système de degré supérieur à 2, on élabore le schéma de
Newton de la manière suivante. Supposons que la solution d'ordre zéro soit le
vecteur $x_i^{(0)}$ et que la solution d'ordre 1 est $x_i^{(1)}$. On a
\[
  f_j\bigl(x^{(1)}\bigr) = f_j\bigl(x^{(0)}\bigr)
  + \sum_{i=1}^{n}\bigl(x_i^{(1)} - x_i^{(0)}\bigr)
    \left.\frac{\partial f_j}{\partial x_i}\right|_{x^{(0)}}
  \qquad (\text{évalué en } x^{(0)}) .
\]
Or $x_i^{(1)}$ étant solution de l'équation, on a $f_j\bigl(x^{(1)}\bigr) = 0$ ;
d'où le système matriciel
\[
  \sum_{i=1}^{n}\bigl(x_i^{(1)} - x_i^{(0)}\bigr)
  \left.\frac{\partial f_j}{\partial x_i}\right|_{x^{(0)}}
  = -f_j\bigl(x^{(0)}\bigr) .
\]
En évaluant ce système matriciel on obtient le vecteur $x_i^{(1)}$. On reprend
le processus et ainsi de suite. À chaque étape on résout le système
matriciel :
\[
  \sum_{i=1}^{n}\bigl(x_i^{(k+1)} - x_i^{(k)}\bigr)
  \left.\frac{\partial f_j}{\partial x_i}\right|_{x^{(k)}}
  = -f_j\bigl(x^{(k)}\bigr) .
\]
Dès que les quantités $x_i^{(k)}$ sont trouvées à la précision choisie, on a
alors l'ensemble des racines du système algébrique non linéaire.
